% GENERATED FILE. Edit canonical lessons, then run npm run generate:books.
% canonical-source-hash: fnv1a64:c46de2fbe7af31a7
% canonical-lessons: JA-C480-tsukuru, JA-C480-deshi, JA-C480-shishou, JA-C480-tema, JA-C480-shiageru

\chapter{Making Things}
\label{ch:ja-making-things}
\hlchaptermodality{480}

\begin{chapteropening}
\textbf{By the end of this chapter:} \emph{I can say who makes things and how long it takes.}

Complete the last lesson of chapter 480: I can say who makes things and how long it takes.
\end{chapteropening}

\section[\texorpdfstring{\ja{つくる} (tsukuru)}{つくる (tsukuru)}]{\ja{つくる} (tsukuru) — to make; to build}

\label{lesson:JA-C480-tsukuru}



\begin{quote}
Before the new one: say the Japanese for serious, then the Japanese for selfish.
\end{quote}

\subsection*{You'll want to know: \ja{つくる}}
\textbf{\ja{つくる}} — \emph{tsukuru} — \textquotedblleft{}to make; to build\textquotedblright{}.

\begin{quote}
\textbf{\ja{ちちは} \ja{いすを} \ja{つくります。}} — \emph{chichi wa isu o tsukurimasu} — \textquotedblleft{}My father makes chairs.\textquotedblright{}
\end{quote}

Polite \textbf{\ja{つくります}}, \emph{tsukurimasu}; a godan verb. Past (polite): \textbf{\ja{ちちが} \ja{たなを} \ja{つくりました。}}, \emph{chichi ga tana o tsukurimashita}.

\begin{cousinweb}
\textbf{\ja{つくる}} (\emph{tsukuru}) is a native Japanese word (wago), inherited rather than borrowed from Chinese; its kanji was attached later as a kun'yomi, and its older origin is not settled (JMdict, KANJIDIC2).
\end{cousinweb}

\begin{grammarlens}[title={What you've built}]
One more word for describing what you lived and felt.
\end{grammarlens}

\subsection*{Your turn}
\begin{itemize}
\raggedright
  \item \emph{Say it:} \emph{tsukuru}
  \item \emph{Say it:} \emph{tsukuru}, once more
  \item \emph{Say it:} \emph{chichi wa isu o tsukurimasu} — My father makes chairs.
  \item \emph{Recall:} say the Japanese for serious, then the Japanese for selfish, then say \emph{tsukuru} again
\end{itemize}

\begin{tcolorbox}[breakable,colback=teal!4,colframe=teal!35!black,title={Before you move on}]
What does \ja{つくる} mean? (\textquotedblleft{}To make; to build\textquotedblright{}.) Say it once more.
\end{tcolorbox}
\section[\texorpdfstring{\ja{でし} (deshi)}{でし (deshi)}]{\ja{でし} (deshi) — an apprentice; a pupil}

\label{lesson:JA-C480-deshi}



\begin{quote}
Before the new one: say the Japanese for selfish, then the Japanese for to make.
\end{quote}

\subsection*{You'll want to know: \ja{でし}}
\textbf{\ja{でし}} — \emph{deshi} — \textquotedblleft{}an apprentice; a pupil\textquotedblright{}.

\begin{quote}
\textbf{\ja{だいくの} \ja{でしは} \ja{まいにち} \ja{はたらきます。}} — \emph{daiku no deshi wa mainichi hatarakimasu} — \textquotedblleft{}The carpenter's apprentice works every day.\textquotedblright{}
\end{quote}

\begin{cousinweb}
\textbf{\ja{でし}} (\emph{deshi}) is Sino-Japanese (kango): two Chinese characters read with their on'yomi, readings borrowed from Chinese. They are \emph{de}, \textquotedblleft{}younger brother\textquotedblright{} (Middle Chinese \emph{dejX}) and \emph{shi}, \textquotedblleft{}child\textquotedblright{} (\emph{tsiX}) (KANJIDIC2; Middle Chinese in Baxter's transcription).
\end{cousinweb}

\begin{grammarlens}[title={What you've built}]
One more word for describing what you lived and felt.
\end{grammarlens}

\subsection*{Your turn}
\begin{itemize}
\raggedright
  \item \emph{Say it:} \emph{deshi}
  \item \emph{Say it:} \emph{deshi}, once more
  \item \emph{Say it:} \emph{daiku no deshi wa mainichi hatarakimasu} — The carpenter's apprentice works every day.
  \item \emph{Recall:} say the Japanese for selfish, then the Japanese for to make, then say \emph{deshi} again
\end{itemize}

\begin{tcolorbox}[breakable,colback=teal!4,colframe=teal!35!black,title={Before you move on}]
What does \ja{でし} mean? (\textquotedblleft{}An apprentice; a pupil\textquotedblright{}.) Say it once more.
\end{tcolorbox}
\section[\texorpdfstring{\ja{ししょう} (shishō)}{ししょう (shishō)}]{\ja{ししょう} (shishō) — a master (of a craft); a teacher}

\label{lesson:JA-C480-shishou}



\begin{quote}
Before the new one: say the Japanese for to make, then the Japanese for an apprentice.
\end{quote}

\subsection*{You'll want to know: \ja{ししょう}}
\textbf{\ja{ししょう}} — \emph{shishō} — \textquotedblleft{}a master (of a craft); a teacher\textquotedblright{}.

\begin{quote}
\textbf{\ja{ししょうは} \ja{でしに} \ja{しごとを} \ja{おしえます。}} — \emph{shishō wa deshi ni shigoto o oshiemasu} — \textquotedblleft{}The master teaches the apprentice the work.\textquotedblright{}
\end{quote}

\begin{cousinweb}
\textbf{\ja{ししょう}} (\emph{shishō}) is Sino-Japanese (kango): two Chinese characters read with their on'yomi, readings borrowed from Chinese. They are \emph{shi}, \textquotedblleft{}expert\textquotedblright{} (Middle Chinese \emph{srij}) and \emph{shō}, \textquotedblleft{}artisan\textquotedblright{} (\emph{dzjangH}) (KANJIDIC2; Middle Chinese in Baxter's transcription).
\end{cousinweb}

\begin{grammarlens}[title={What you've built}]
One more word for describing what you lived and felt.
\end{grammarlens}

\subsection*{Your turn}
\begin{itemize}
\raggedright
  \item \emph{Say it:} \emph{shishō}
  \item \emph{Say it:} \emph{shishō}, once more
  \item \emph{Say it:} \emph{shishō wa deshi ni shigoto o oshiemasu} — The master teaches the apprentice the work.
  \item \emph{Recall:} say the Japanese for to make, then the Japanese for an apprentice, then say \emph{shishō} again
\end{itemize}

\begin{tcolorbox}[breakable,colback=teal!4,colframe=teal!35!black,title={Before you move on}]
What does \ja{ししょう} mean? (\textquotedblleft{}A master (of a craft); a teacher\textquotedblright{}.) Say it once more.
\end{tcolorbox}
\section[\texorpdfstring{\ja{てま} (tema)}{てま (tema)}]{\ja{てま} (tema) — time and effort; trouble}

\label{lesson:JA-C480-tema}



\begin{quote}
Before the new one: say the Japanese for an apprentice, then the Japanese for a master (of a craft).
\end{quote}

\subsection*{You'll want to know: \ja{てま}}
\textbf{\ja{てま}} — \emph{tema} — \textquotedblleft{}time and effort; trouble\textquotedblright{}.

\begin{quote}
\textbf{\ja{この} \ja{しごとは} \ja{てまが} \ja{かかります。}} — \emph{kono shigoto wa tema ga kakarimasu} — \textquotedblleft{}This work takes a lot of time and effort.\textquotedblright{}
\end{quote}

\begin{cousinweb}
\textbf{\ja{てま}} (\emph{tema}) is a native Japanese compound (wago), inherited rather than borrowed from Chinese: built from \emph{te}, \textquotedblleft{}hand\textquotedblright{} and \emph{ma}, \textquotedblleft{}interval\textquotedblright{} (JMdict, KANJIDIC2).
\end{cousinweb}

\begin{grammarlens}[title={What you've built}]
One more word for describing what you lived and felt.
\end{grammarlens}

\subsection*{Your turn}
\begin{itemize}
\raggedright
  \item \emph{Say it:} \emph{tema}
  \item \emph{Say it:} \emph{tema}, once more
  \item \emph{Say it:} \emph{kono shigoto wa tema ga kakarimasu} — This work takes a lot of time and effort.
  \item \emph{Recall:} say the Japanese for an apprentice, then the Japanese for a master (of a craft), then say \emph{tema} again
\end{itemize}

\begin{tcolorbox}[breakable,colback=teal!4,colframe=teal!35!black,title={Before you move on}]
What does \ja{てま} mean? (\textquotedblleft{}Time and effort; trouble\textquotedblright{}.) Say it once more.
\end{tcolorbox}
\section[\texorpdfstring{\ja{しあげる} (shiageru)}{しあげる (shiageru)}]{\ja{しあげる} (shiageru) — to finish off; to complete (a piece of work)}

\label{lesson:JA-C480-shiageru}



\begin{quote}
Before the new one: say the Japanese for a master (of a craft), then the Japanese for time and effort.
\end{quote}

\subsection*{You'll want to know: \ja{しあげる}}
\textbf{\ja{しあげる}} — \emph{shiageru} — \textquotedblleft{}to finish off; to complete (a piece of work)\textquotedblright{}.

\begin{quote}
\textbf{\ja{あしたまでに} \ja{いすを} \ja{しあげます。}} — \emph{ashita made ni isu o shiagemasu} — \textquotedblleft{}I will finish the chair by tomorrow.\textquotedblright{}
\end{quote}

Polite \textbf{\ja{しあげます}}, \emph{shiagemasu}; an ichidan verb. Past (polite): \textbf{\ja{だいくが} \ja{いすを} \ja{しあげました。}}, \emph{daiku ga isu o shiagemashita}.

\begin{cousinweb}
\textbf{\ja{しあげる}} (\emph{shiageru}) is a native Japanese word (wago), inherited rather than borrowed from Chinese; its kanji was attached later as a kun'yomi, and its older origin is not settled (JMdict, KANJIDIC2).
\end{cousinweb}

\begin{grammarlens}[title={What you've built}]
One more word for describing what you lived and felt.
\end{grammarlens}

\subsection*{Your turn}
\begin{itemize}
\raggedright
  \item \emph{Say it:} \emph{shiageru}
  \item \emph{Say it:} \emph{shiageru}, once more
  \item \emph{Say it:} \emph{ashita made ni isu o shiagemasu} — I will finish the chair by tomorrow.
  \item \emph{Recall:} say the Japanese for a master (of a craft), then the Japanese for time and effort, then say \emph{shiageru} again
\end{itemize}

\begin{tcolorbox}[breakable,colback=teal!4,colframe=teal!35!black,title={Before you move on}]
What does \ja{しあげる} mean? (\textquotedblleft{}To finish off; to complete (a piece of work)\textquotedblright{}.) Say it once more.
\end{tcolorbox}
